\section{Design problems and sensitivities}\label{sec:design}

\subsection{Orientation field and the tow-curvature constraint}\label{sec:manuf}
The fiber angle is a smooth field interpolated from $M$ support values $\hat{\bm\theta}$
by normalized compactly supported Wendland $C^2$ radial basis functions \cite{wendland1995},
\begin{equation}
\theta(\bm x)=\sum_{i=1}^M\Phi_i(\bm x)\hat\theta_i,\qquad
\Phi_i=\frac{\phi(\|\bm x-\bm x_i\|/r_s)}{\sum_j\phi(\|\bm x-\bm x_j\|/r_s)},\qquad
\phi(q)=(1-q)_+^4(4q+1),
\label{eq:csrbf}
\end{equation}
on a regular support grid of spacing $h_s$ with $r_s=1.5h_s$. The baseline is
$h_s=\SI{50}{mm}$ ($M=20$); Sec.~\ref{sec:limits} refines it. The spatial derivatives
$\bm\Phi_x,\bm\Phi_y$ are analytic. A tow laid along $\bm n=(\cos\theta,\sin\theta)$
turns at the rate
\begin{equation}
\kappa_{\mathrm{tow}}(\bm x)=(\bm n\!\cdot\!\nabla)\theta=\cos\theta\,\theta_{,x}+\sin\theta\,\theta_{,y},
\label{eq:towcurv}
\end{equation}
so its steering radius is $1/|\kappa_{\mathrm{tow}}|$. AFP and continuous-fiber printing require
$|\kappa_{\mathrm{tow}}|\le\zall=1/R_{\min}$; this pointwise form follows \cite{wong2026}. We impose this at every quadrature point $e$
as $c_e=1-\kappa_{\mathrm{tow},e}^2/\zall^2\ge0$, with the analytic Jacobian
\begin{equation}
\frac{\partial\kappa_{\mathrm{tow},e}}{\partial\hat{\bm\theta}}=
\big(-\sin\theta_e\,\bm\Phi_{x,e}\hat{\bm\theta}+\cos\theta_e\,\bm\Phi_{y,e}\hat{\bm\theta}\big)\bm\Phi_e
+\cos\theta_e\,\bm\Phi_{x,e}+\sin\theta_e\,\bm\Phi_{y,e}.
\end{equation}
The baseline limit is $R_{\min}=\SI{20}{mm}$. The constraint is necessary for laying the field with
tows but not sufficient: tow spreading and the resulting gaps or overlaps are treated separately
(Sec.~\ref{sec:asbuilt}).

\subsection{The design problems}\label{sec:problems}
\paragraph{Steered plate.} For a target frequency $\omega_\star$,
\begin{equation}
\min_{\hat{\bm\theta}}\;\frac{J(\hat{\bm\theta})}{J_{\UD}}\quad\text{s.t.}\quad
\frac{\omega_1(\hat{\bm\theta})}{\omega_\star}-1=0,\qquad c_e(\hat{\bm\theta})\ge0\ \ \forall e,
\label{eq:designprob}
\end{equation}
where $J$ and $\omega_1$ are those of the adapted-mode, relaxed model
(Sec.~\ref{sec:relaxed}). The frequency is an \emph{equality}, so a gain can never come
from a stiffer or softer panel. The mass is fixed by construction. We solve
\eqref{eq:designprob} by sequential quadratic programming (SLSQP \cite{kraft1988}), which satisfies both
constraints to $10^{-4}$, from six starts (uniform $0^\circ,20^\circ,45^\circ,70^\circ$ and
two random fields), and keep the best feasible result.

\paragraph{Uniform stacks.} We enumerate every symmetric eight-ply laminate
$[\theta_4/\theta_3/\theta_2/\theta_1]_s$ with $\theta_i\in\{-75^\circ,-60^\circ,\dots,90^\circ\}$:
$12^4=\nSnum$ stacks, each evaluated exactly on the same model. A uniform laminate
needs no optimizer, because its operators are linear combinations of precomputed moment
tensors (e.g.\ $\bm K_b=\sum_{ab}D_{ab}\bm I^K_{ab}$).

\paragraph{Steered stack.} The best uniform stack at $\omega_\star$ is rotated pointwise,
$\bm C(\bm x)=\bm T(\theta(\bm x))\bm C_0\bm T(\theta(\bm x))^{\!\top}$ for $\bm C\in\{\bm A,\bm
D\}$, and \eqref{eq:designprob} is solved for $\theta(\bm x)$.

\paragraph{Beam.} $\theta(x)$ is mirror symmetric about mid-span and piecewise linear on
13 knots over the half span. The tow limit is the slope bound
$|\theta'|\le\kappa_{\mathrm{tow}}$, and the frequency is held as
$\omega_1\ge f_{\mathrm{hold}}\,\omega_{1,\UD}$. A uniform-angle beam cannot change $J$
(Corollary~\ref{cor:beaminv}) and every steered beam is \emph{softer} than UD, so an
equality at $f_{\mathrm{hold}}=1$ admits only UD. We report
$f_{\mathrm{hold}}\in\{0.98,0.90\}$ and $\kappa_{\mathrm{tow}}\in\{900,1800,\infty\}$\,deg/m.
The optimizer works on the condensed beam evaluator (adapted mode, exact axial
condensation), and every design is \emph{confirmed} on the full FE beam by its cubic SSM.

\subsection{Sensitivities on the adapted mode}\label{sec:sens}
Let a perturbation $\delta\theta(\bm x)$ change the stiffnesses by
$\delta\bm D=\bm D'\delta\theta$ and $\delta\bm A=\bm A'\delta\theta$, with
$\bm C'=\bm T'\bm C_0\bm T^{\!\top}+\bm T\bm C_0\bm T'^{\!\top}$. Let
$(\lambda_j,\bm v_j)$ be the Ritz pairs of the bending block ($\lambda_1=\omega_1^2$),
with curvature fields $\bm\kappa_j$. The mass does not depend on $\theta$.

\emph{Frequency.} The Rayleigh quotient is stationary (Fox--Kapoor), so
\begin{equation}
\delta\omega_1^2=\int\bm\kappa_1^{\!\top}\delta\bm D\,\bm\kappa_1 .
\label{eq:dom}
\end{equation}

\emph{Mode shape.} The eigenvector derivative is
\begin{equation}
\delta\bm v_1=\sum_{j\ge2}\frac{\bm v_j^{\!\top}\delta\bm K_b\bm v_1}{\lambda_1-\lambda_j}\bm v_j,
\qquad \bm v_j^{\!\top}\delta\bm K_b\bm v_1=\int\bm\kappa_j^{\!\top}\delta\bm D\,\bm\kappa_1 .
\end{equation}

\emph{Hardening.} By Proposition~\ref{prop:relaxed}, $\gamma$ is a minimum over the
in-plane state, so its total derivative needs \emph{no in-plane adjoint}: only the
explicit $\bm A$-dependence and the dependence on the mode (through $\bm m$) remain.
With $\partial\gamma/\partial\bm m=\bm A\bm\varepsilon_r$ and
$\partial\bm m/\partial v_k=[2w_{1,x}\phi_{k,x},\,2w_{1,y}\phi_{k,y},\,
2(w_{1,x}\phi_{k,y}+w_{1,y}\phi_{k,x})]$,
\begin{equation}
\delta\gamma=\int\Big[\tfrac12\bm\varepsilon_r^{\!\top}\delta\bm A\,\bm\varepsilon_r
+\bm\kappa_\beta^{\!\top}\delta\bm D\,\bm\kappa_1\Big],\qquad
\bm\kappa_\beta=\sum_{j\ge2}\frac{\bm v_j^{\!\top}\partial_{\bm v}\gamma}{\lambda_1-\lambda_j}\bm\kappa_j,\quad
\partial_{v_k}\gamma=\int(\bm A\bm\varepsilon_r)^{\!\top}\partial_{v_k}\bm m .
\label{eq:dgamma}
\end{equation}
The \emph{adjoint curvature field} $\bm\kappa_\beta$ collapses every eigenvector term into
one field.

\emph{Mid-plate value.} In the same way,
$\delta w_c=\int\bm\kappa_{\mathrm{mid}}^{\!\top}\delta\bm D\,\bm\kappa_1$, with
$\bm\kappa_{\mathrm{mid}}=\sum_{j\ge2}\frac{w_c(\bm v_j)}{\lambda_1-\lambda_j}\bm\kappa_j$.

\emph{Detuning.} Since $J=\gamma/(\omega_1^2w_c^2)$,
\begin{equation}
\delta J=\frac{\delta\gamma}{\omega_1^2w_c^2}-\frac{\gamma\,\delta\omega_1^2}{\omega_1^4w_c^2}
-\frac{2\gamma\,\delta w_c}{\omega_1^2w_c^3},\qquad
\frac{\partial J}{\partial\hat{\bm\theta}}=\bm\Phi^{\!\top}\Big[w_q\frac{\partial J}{\partial\theta_q}\Big]_q ,
\label{eq:dJ}
\end{equation}
where $w_q$ are the quadrature weights. The whole gradient costs one eigensolve of
dimension $N_w$, one in-plane solve of dimension $2N_s$, and pointwise contractions. It
does not depend on $M$, so the support grid can be refined freely.
Equation~\eqref{eq:dJ} is the local optimality condition of
Proposition~\ref{prop:pointwise}, plus the mode-shape corrections carried by
$\bm\kappa_\beta$ and $\bm\kappa_{\mathrm{mid}}$.

\subsection{Verification of the evaluator and its gradient}\label{sec:verify}
\begin{itemize}[itemsep=1pt]
\item \textbf{Gradient.} At the published design, central finite differences (steps $10^{-4}$--$10^{-6}$)
  along $\nGradDirs$ directions (five random, five largest-gradient coordinates) agree with the analytic
  derivatives of $J$, $\omega_1$, $\gamma$ and the tow curvature to better than $\nGradErr$ relative
  (\texttt{r1\_\allowbreak verification.py}).
\item \textbf{Constraints on an independent grid.} The tow curvature of the final field, evaluated
  analytically on a 1\,mm grid that shares no point with the optimizer's quadrature, gives a minimum
  radius of $\nRminFine$\,mm (limit 20\,mm). The frequency equality is re-checked on finer bases and on
  the FE model in Sec.~\ref{sec:confirm}.
\item \textbf{Multi-start.} Table~\ref{tab:ms} lists every start of the held-frequency solve.
\item \textbf{Isotropy.} With $\bm Q$ isotropic, $\partial J/\partial\theta$ vanishes to
  $10^{-10}$ relative (Corollary~\ref{cor:iso}).
\item \textbf{Proposition~\ref{prop:relaxed}.} $\int e_G=\gamma$ and $\int e_K=\omega_1^2$
  hold to $10^{-10}$.
\item \textbf{Rotation.} The rotation \eqref{eq:rot} reproduces the classical $\bar{\bm Q}$
  to machine precision, and Proposition~\ref{prop:fourier} holds to $10^{-9}$ for random
  strains.
\item \textbf{Uniform path.} The fast uniform-laminate path equals the general one to
  $10^{-13}$.
\item \textbf{Reference values.} For the UD plate the evaluator gives $f_1=\nUDf$\,Hz and a
  $\nUDdet\%$ shift at $a/h=0.4$, in agreement with the $\nFEtabUDdet\%$ of the assumed-modes
  model of the finite-element comparison (Appendix~\ref{app:fe}).
\end{itemize}
Basis convergence is in Table~\ref{tab:conv}.

\begin{table}[!tbp]
\centering\footnotesize
\caption{Multi-start diagnostics of the held-frequency plate design (all SLSQP runs terminated
normally; residuals at termination; generated by \texttt{r1\_\allowbreak verification.py}). The reported
design is the best start; $\nMSgood$ of six starts reach within $\nMSspread\%$ of it, and two stall in
poorer local minima.}
\label{tab:ms}
\IfFileExists{ms_table.tex}{\begin{tabular}{@{}lcccc@{}}
\toprule
start & $J/J_{\UD}$ & $\omega_1/\omega_\star-1$ & $\min_e c_e$ & iterations\\
\midrule
uniform 0 & $0.6508$ & $3\times10^{-13}$ & $-2\times10^{-14}$ & 22\\
uniform 20 & $0.3659$ & $3\times10^{-11}$ & $8\times10^{-12}$ & 54\\
uniform 45 & $0.3813$ & $9\times10^{-14}$ & $-4\times10^{-15}$ & 54\\
uniform 70 & $0.5711$ & $3\times10^{-13}$ & $7\times10^{-14}$ & 44\\
random 1 & $0.3675$ & $3\times10^{-11}$ & $1\times10^{-12}$ & 72\\
random 2 & $0.3675$ & $2\times10^{-11}$ & $4\times10^{-12}$ & 79\\
\bottomrule
\end{tabular}
}{\textit{table pending}}
\end{table} All checks are in
\texttt{tests/test\_relaxed\_plate.py}.

\begin{table}[!tbp]
\centering\footnotesize
\caption{Convergence of the evaluator on the reported held-frequency steered design
(generated by \texttt{conv\_table.py}). $J/J_{\UD}$ and the relaxed share are converged on
the bold basis; three in-plane modes per direction under-resolve the steered relaxation.
The steered frequency is a Ritz upper bound that falls by $\nConvFreqDrop\%$ as the transverse basis
grows (the UD frequency does not change), so ``held frequency'' holds to about $\nConvFreqDrop\%$
across bases and $\nHeldTol\%$ against the independent FE model.
Bold: the basis used throughout.}
\label{tab:conv}
\IfFileExists{conv_table.tex}{\begin{tabular}{@{}ccccccc@{}}
\toprule
transverse $m,n\le$ & in-plane $i,j\le$ & $f_{1,\UD}$ (Hz) & $f_1$ (Hz) & $J/J_{\UD}$ & relaxed share (UD) & relaxed share (steered)\\ \midrule
8 & 3 & 256.3 & 256.3 & 0.523 & 33\% & 31\%\\
8 & 4 & 256.3 & 256.3 & 0.420 & 33\% & 44\%\\
8 & 5 & 256.3 & 256.3 & 0.375 & 33\% & 50\%\\
\bf 8 & \bf 6 & \bf 256.3 & \bf 256.3 & \bf 0.366 & \bf 33\% & \bf 51\%\\
8 & 8 & 256.3 & 256.3 & 0.360 & 33\% & 52\%\\
12 & 6 & 256.3 & 254.6 & 0.370 & 33\% & 51\%\\
12 & 8 & 256.3 & 254.6 & 0.364 & 33\% & 52\%\\
16 & 8 & 256.3 & 253.9 & 0.366 & 33\% & 51\%\\
\bottomrule
\end{tabular}
}{\textit{table pending}}
\end{table}
